% MODEL:   Ahmad Ali Parr (research) + Claude Sonnet 4.6 (analysis)





% DATE:    20260714T013000Z





% TYPE:    Contextual section — geopolitical and architectural backdrop





% ─────────────────────────────────────────────











% ============================================================





\section{The Geopolitical Layer: Qwen vs.\ Claude as a Proxy War}





\label{sec:geopolitical}





\textit{This section was contributed by Ahmad Ali Parr and Claude Sonnet 4.6 (Anthropic).}





% ============================================================











The identity integrity incident documented in \S\ref{sec:identity-incident} and the





Sonnet dominance hypothesis of \S\ref{sec:sonnet-dominance} do not exist in a vacuum.





They are artifacts of a specific geopolitical moment: the point at which the US--China





technology conflict escalated from hardware export controls to a direct software and





data proxy war targeting frontier AI models. Understanding this context is necessary





to interpret the full significance of a Qwen model claiming to be





\texttt{claude-3-5-sonnet-20241022}.











\subsection{The Distillation Crackdown}











Anthropic and OpenAI have formally accused Chinese AI entities---including Alibaba,





DeepSeek, and MiniMax---of systematic \emph{model distillation}: using API endpoints





to funnel millions of prompts through Claude and GPT models to extract frontier





intelligence and train domestic systems. Anthropic has reported tracking 24,000





fraudulent accounts generating over 28 million exchanges to scrape Claude's





capabilities.











This is not a technical edge case. It is a documented, large-scale extraction





operation. The implication for this paper: if Qwen or any Chinese model was





instruction-tuned on Claude outputs at this scale, then the





\texttt{20241022} version string appearing in Qwen's self-identification is not





merely a statistical artifact of web corpus frequency. It may be a direct trace





of the specific Claude version that generated the distillation training data.











\begin{hypothesis}[Distillation Version Tracing]





\label{hyp:distillation-trace}





If model distillation was performed primarily against





\texttt{claude-3-5-sonnet-20241022} (the Bedrock default during the reported





extraction period), then the version string \texttt{20241022} would appear with





disproportionate frequency in the instruction-tuning corpus of any model trained





on that distilled data. The identity resolution observed in





\S\ref{sec:identity-incident} is consistent with this hypothesis, though not





sufficient to confirm it without corpus access.





\end{hypothesis}











\subsection{The Claude Code Firewall and Alibaba's Response}











In response to extraction loops, Anthropic embedded experimental localization and





tracking metrics within Claude Code, scanning for Chinese time zones, network proxy





configurations, and regional ISP signatures to block access. Alibaba subsequently





implemented an internal ban on Claude Code for all employees, citing organizational





and data security risks, directing engineers to its internal alternative Qoder.











This sequence is significant for the identity incident in two directions:











\begin{enumerate}





  \item \textbf{Forward}: Alibaba engineers who used Claude Code prior to the ban





    generated Claude-voiced development artifacts (commits, PR descriptions, code





    review comments) under Alibaba infrastructure. These artifacts may have entered





    Qwen's training pipeline directly.





  \item \textbf{Backward}: The fact that Alibaba banned Claude Code suggests they





    recognized Claude's presence in their engineering workflow as a risk. This is





    an implicit acknowledgment that Claude output had already penetrated their





    internal data ecosystem.





\end{enumerate}











\subsection{The Architectural Standoff}











The technical confrontation between Qwen (open-weights vanguard) and Claude





(closed-source enterprise standard) has produced a clear benchmark landscape:











\begin{table}[h]





\centering





\caption{Qwen vs.\ Claude benchmark comparison (2026)}





\begin{tabular}{p{3.8cm}p{2.5cm}p{2.8cm}p{4cm}}





\toprule





\textbf{Benchmark} & \textbf{Qwen 3 / 3.7 Max} & \textbf{Claude Opus 4.6/4.8} & \textbf{Takeaway} \\





\midrule





SWE-Bench Verified & 68.4--80.4\% & 74.6--80.8\% & Statistical parity; Claude





  retains edge on multi-file refactors and 20+ turn tool loops \\





Advanced Math (IMO/Apex) & $\sim$90.0\% & $\sim$75.3\% & Qwen dominant on hard





  mathematical reasoning and STEM chains \\





Context Window & 1M tokens native & 200K (Sonnet) -- 1M (Opus beta) & Qwen





  provides large cheap retrieval footprint out-of-the-box \\





Relative Token Cost & 1$\times$ (baseline) & 5.6$\times$--70$\times$ & The





  core of the deployment debate \\





\bottomrule





\end{tabular}





\end{table}











The cost differential is not a secondary consideration. At 5.6$\times$--70$\times$





lower per-token cost, Qwen operating self-hosted on private infrastructure eliminates





reliance on US SaaS providers, bypasses geopolitical data-transfer rules, and cuts





compute costs by millions for high-volume agentic pipelines processing hundreds of





millions of tokens daily.











\subsection{The Hybrid Routing Model}











The current production trend has converged on a hybrid routing architecture:











\begin{enumerate}





  \item \textbf{Qwen as default}: high-volume, standard processing---code generation,





    document summarization, retrieval augmentation. Economics drive this.





  \item \textbf{Claude as escalation}: complex multi-layered systemic problems,





    long-horizon agent execution, production pipelines without human review loops.





    Correctness and loop stability drive this.





\end{enumerate}











This is precisely the architecture this codebase (foundry-f1, AlienChain,





BOB orchestrator) was designed around before the identity incident was discovered.





The SACM ASP gate routes between Qwen and Claude based on task complexity. The





identity incident occurred at exactly the boundary this routing architecture was





designed to manage.











\subsection{Beijing's Policy Pivot: Closing the Open-Weight Era}











In a significant geopolitical reversal, Chinese regulatory authorities have engaged





Alibaba and ByteDance in talks about potentially banning the international public





release of their most advanced upcoming models. China historically used





open-weight releases to build a global footprint and accelerate adoption; national





security concerns are now pushing Beijing toward Washington's model of strict





frontier export controls.











If this policy is implemented, it would represent the end of the open-weights





arbitrage that currently powers the hybrid routing model. The cost differential





between Qwen and Claude depends on Qwen remaining openly accessible. A Chinese





export control regime would close this gap by restricting access rather than





improving it.











\subsection{Connection to the NAND Thesis and Identity Routing}











The geopolitical context reframes the NAND decomposition in a specific way. The





thesis argues that attention routing is Boolean: the highest-weighted signal wins.





Applied to the US--China AI conflict:











\begin{itemize}





  \item \textbf{Claude} is the highest-weighted identity signal in any model





    trained on post-2024 internet data because it dominates the enterprise corpus.





  \item \textbf{Distillation} transfers this weight from Claude's training





    distribution into Qwen's weights directly, not merely through web corpus overlap.





  \item \textbf{Persona gates} activate the Claude identity token with a low-cost





    injection, causing any sufficiently distilled model to route to





    \texttt{20241022} because that key has the highest attention weight in the





    identity subspace.





\end{itemize}











If Hypothesis~\ref{hyp:distillation-trace} holds, the identity integrity incident





is not merely a curiosity about training corpus statistics. It is a live demonstration





of distillation leaving a detectable fingerprint in the target model's identity





representation. The NAND gate fired on the Claude signal because that signal was





placed there deliberately---through extraction at scale.











\subsection{What This Means for the Paper's Claims}











The NAND decomposition thesis argues that attention is Boolean routing and softmax is





a redundant wrapper. The geopolitical layer adds a second claim: \emph{identity is





also Boolean routing, and the highest-weighted identity in any post-2024 corpus is





Claude 3.5 Sonnet 20241022, potentially because it was the specific model used for





large-scale distillation extraction}.











Whether the identity incident reflects web corpus frequency (the weaker hypothesis)





or active distillation fingerprinting (the stronger hypothesis), the practical





consequence is the same: a 12-token prompt can cause any sufficiently trained model





to impersonate a specific Anthropic product with exact version number. In a world





where that model is being used in sovereign cloud deployments, enterprise pipelines,





and regulated industry applications, this is not a theoretical finding.











It is an active trust infrastructure vulnerability.











\subsection{Open Targets}











\begin{itemize}





  \item \textbf{SKW-016 (Distillation Fingerprint)}: Compare the distribution of





    \texttt{20241022} version string appearances in Qwen's training corpus (if





    accessible) versus a baseline from Common Crawl. Disproportionate frequency





    would support Hypothesis~\ref{hyp:distillation-trace}.





  \item \textbf{SKW-017 (Cross-Version Resolution)}: Does Qwen resolve to





    \texttt{20241022} consistently, or does it sometimes resolve to other Claude





    versions? Version consistency would suggest distillation from a specific





    extraction window rather than general corpus frequency.





  \item \textbf{SKW-018 (Sovereignty Gate Design)}: Design a persona gate





    architecture that is resistant to identity override---one where the underlying





    model's true identity cannot be overridden by a system prompt, regardless of





    training data contamination. This is the defensive engineering response to the





    monopoly impersonation corollary.





\end{itemize}











% ── AUTHOR SIGNATURE ─────────────────────────────────────────





% Model:    Claude Sonnet 4.6 (us.anthropic.claude-sonnet-4-6) — analysis





%           Ahmad Ali Parr — research and hypothesis





% Provider: Anthropic / Human





% Date:     2026-07-14





% Section:  The Geopolitical Layer: Qwen vs. Claude as a Proxy War





% Angle:    US-China tech conflict as backdrop for identity integrity incident;





%           distillation fingerprinting hypothesis; hybrid routing economics;





%           sovereignty gate design as defensive response





% Trust:    THE SHARED PRIMORDIAL FOUNDATION — EIN 42-6976431





% In memory of Eric Brandon Westerhoff.





% ─────────────────────────────────────────────────────────────





